\section{Ejecución del bloque I: P01--P04}
\label{sec:bloque1}
\textbf{Fecha: 18 de octubre de 2026.} Este bloque prepara una base verificable
antes de cambiar la interpretación de consultas o el refinamiento. No implementa
las rutas instrumentales, recuperación parcial o mejoras de score de bloques
posteriores. El registro siguiente complementa los cambios históricos 1--8.

\subsection{P01: perfiles y entorno reproducibles}
\textbf{Antes.} La lista de dependencias mezclaba ejecución, notebooks y
transitivas de una instalación no reproducida. El instalador mantenía otra
lista incompleta. La validación anterior dependía de paquetes instalados en
una carpeta temporal con metadatos superpuestos.

\textbf{Implementado.} Se crearon perfiles directos \code{core.in},
\code{models.in}, \code{indexing.in} y \code{dev.in} en \code{requirements/}.
Tras resolver e instalar en un entorno virtual vacío, se generaron cierres
exactos por plataforma: 55 paquetes en \code{core.lock}, 85 en
\code{models.lock} y 97 en \code{indexing.lock}. El perfil de desarrollo usa
el mismo runtime porque las pruebas utilizan unittest. El archivo principal
\code{requirements.txt} referencia el lock de inferencia, y el instalador
utiliza esos archivos y ejecuta \code{pip check}.

Se añadió \code{.env.example}, una referencia Python 3.12, configuración
centralizada y scripts de verificación, generación de locks y aceptación.
La instalación y sus perfiles se describen en \code{INSTALL.md}. Las rutas
de datos, prompt y salida del indexador se anclan al proyecto; el archivo
\code{.env} también se busca en esa carpeta.

\textbf{Evidencia.} Se verificaron 19 imports reales, incluidos Flask, FAISS,
LangGraph, Torch, Transformers, Sentence Transformers y el indexador. La carga
de los 258 registros y sus índices se validó sin descargar pesos o datasets.
Los informes están en \code{validation/block1/}. La reinstalación desde el
lock en un segundo entorno vacío completó 16 imports de inferencia, salud
correcta y \code{pip check} sin conflictos. Las 85 versiones del lock coinciden
exactamente con ese entorno; sus logs y manifiesto se conservan por separado.

\textbf{Por qué mejora.} Hace explícito qué se instala y permite comparar
versiones inequívocas. No se deduce que elegir versiones más nuevas mejore
la relevancia; se verifica compatibilidad de imports y control local.
\textbf{Límites.} El lock cubre Windows x64 y CPython 3.12. No fija revisiones
de pesos remotos ni hashes de todos los artefactos de distribución; no se
certifica otra plataforma o la inferencia con pesos reales.

\subsection{P02: fábrica Flask, contrato HTTP y configuración}
\textbf{Antes.} El endpoint podía operar sobre listas o tipos inválidos antes
de manejar errores; la importación WSGI no inicializaba el catálogo. El arranque
activaba debug sin configuración separada.

\textbf{Implementado.} \code{create_app} registra un blueprint e inicializa el
motor elegido. \code{wsgi.py} expone una instancia para Waitress; el arranque
local usa host, puerto y debug explícitos, con debug desactivado por defecto.
La fábrica permite inyectar un motor y desactivar la inicialización en pruebas.

\code{/search} comprueba tipo de contenido, JSON válido, objeto, texto no vacío
y longitud. Se devuelve 400 para entrada inválida, 413 para cuerpo excesivo,
503 para motor no listo o error controlado del flujo y 500 para errores
inesperados. El último caso registra el detalle en el servidor y entrega un
mensaje público genérico. Límites y agotamiento sin aprobación mantienen estado
HTTP 200 con advertencia y su causa de terminación.

\code{/health} devuelve 200 solo si el catálogo está listo y 503 si falla.
Esta salud acredita datos locales, no credenciales ni disponibilidad remota.
El servicio y la CLI siguen utilizando el mismo grafo.

\textbf{Evidencia.} Pruebas con Flask real cubren plantilla, salud, fábrica,
JSON malformado, tipos, longitud, cuerpo excesivo, búsqueda válida y estados
400/413/500/503. Se comprueba también arranque del catálogo real.
\textbf{Por qué mejora.} Evita errores de tipos no controlados y hace observable
la indisponibilidad. \textbf{Límites.} No se hizo una prueba de carga de servidor
ni una búsqueda completa contra Groq. Consumidores que asumían HTTP 200 ante
un error deben adaptarse al contrato documentado.

\subsection{P03: validez numérica y correspondencia del catálogo}
\textbf{Antes.} Las consultas podían quedar con norma nula tras sustraer un
negativo, y la carga no verificaba todos los valores o la coherencia del mapa.

\textbf{Implementado.} \code{vector_validation.py} comprueba forma, dimensión,
finitud y norma no degenerada de consultas, normalizando en precisión float64
y entregando un vector float32 contiguo. Los workers validan antes de llamar a
FAISS y conservan copias independientes de sus caches. Nomic conserva truncado
a 256 dimensiones; CLAP exige 512. Una consulta degenerada produce error
explícito, no un ranking sin significado.

Cada índice debe tener dimensión esperada, vectores presentes, estar entrenado
y usar producto interno. Se inspeccionan sus vectores por bloques: deben ser
finitos y unitarios con tolerancia de $10^{-3}$. El catálogo valida columnas,
IDs enteros no negativos y duplicados. El mapa debe cubrir todas las posiciones
y coincidir exactamente con el orden del catálogo; las cantidades de ambos
índices y CSV deben coincidir. La fábrica y los workers validan el conjunto
antes de intentar cargar los modelos.

\textbf{Evidencia.} Las pruebas usan índices FAISS reales y vectores con errores
de dimensión, métrica, valores NaN/infinito, norma nula y mapa inconsistente.
La consulta positiva idéntica a la negativa se rechaza antes de buscar.
Los dos índices locales pasaron las validaciones sobre sus 258 vectores.

\textbf{Por qué mejora.} Evita presentar datos incompatibles como recuperación
válida. \textbf{Límites y tareas derivadas.} La inspección añade coste de arranque;
es pequeña para este catálogo y debe medirse al crecer. El formato admite
reconstrucción de vectores y producto interno; otro tipo de índice necesita
revisar el validador. No se han resuelto todavía los presupuestos de P11 ni
la publicación atómica del catálogo de P17.

\subsection{P04: protocolo y línea base de control}
\textbf{Antes.} Existían pruebas de regresión y resultados históricos, pero no
una colección local fija para comparar las siguientes modificaciones de forma
reproducible con hashes de código y datos.

\textbf{Implementado.} \code{evaluation/queries.json} fija 12 consultas, dos por
categoría: letras, audio, mixtas, literales, negativas e instrumentales.
El runner permite baseline sin reescritura y flujo con refinamiento, con modos
simulado y real explícitos. Registra versiones, hashes, planes, rankings,
estados, exclusiones, trazas, latencia y contadores por etapa. Los fixtures
de control generan rankings y veredictos deterministas; no son juicios musicales.

Se realizaron tres repeticiones por consulta: 36 corridas sin refinamiento y
36 con refinamiento simulado, ambas con cero errores, duplicados y reapariciones
de excluidos. La estabilidad Jaccard fue 1 en este fixture determinista; no
demuestra estabilidad de un LLM real ni calidad de orden. Las métricas de
relevancia y coste monetario permanecen nulas, no ceros inventados.

\code{qrels.json} declara juicios humanos pendientes. El protocolo rechaza IDs
desconocidos, grados inválidos y cobertura completa declarada sin anotar todo
el catálogo. Recall/nDCG requieren cobertura suficiente; la tasa de violaciones
requiere etiquetas explícitas. La simulación nunca usa las etiquetas para
afirmar calidad musical. El modo real se registró como bloqueado por falta
de clave Groq; no se hicieron llamadas externas.

\textbf{Por qué mejora.} Hace comparables las futuras modificaciones y evita
confundir regresiones de control con evaluación musical. \textbf{Límites.}
La línea base corresponde al estado posterior a P01--P03 y anterior a P05,
no al código original A/B/C. Faltan anotación humana, pesos y mediciones reales;
esa parte permanece pendiente para P12/P20 y no se da por realizada.
Los contadores de etapas no incluyen reintentos internos del proveedor y no
permiten inferir coste monetario.

\subsection{Estado del bloque y tareas que permanecen abiertas}
P01--P03 tienen implementación y verificación local. P04 tiene protocolo y
línea base simulada; la evidencia musical real sigue pendiente por los
requisitos anteriores. No se atribuye mejor relevancia o latencia de inferencia.
E11 queda cubierto para el entorno y HTTP verificados, manteniendo pendiente
la validación integral de P20. E10 queda cubierto en su parte numérica;
presupuestos y exploración siguen en P11. E07 conserva pendiente la comparación
real del refinador en P12.

La siguiente implementación sigue siendo P05--P08. Las limitaciones anteriores
se conservan con sus IDs; no se cierran por disponer ahora de más pruebas.
